---
format: yohaku/0.2
target: text
macro: '[[macros/raw]]'
---

## 正文
````text
"""Yohaku-only provider selection and credential-isolating Responses relay.

Never edits shared Codex settings. The relay forwards only to DeepSeek HTTPS;
its evidence contains metadata, never headers, prompts, output text, or secrets.
"""
import argparse
import getpass
import hmac
import http.client
from http.server import BaseHTTPRequestHandler, HTTPServer
import json
import os
from pathlib import Path
import secrets
import stat
import sys
import tempfile
import threading
import warnings
import urllib.error
from yohaku_provider_transport import Recovery, RecoveryCancelled, ObservedHTTPS, retryable
from yohaku_provider_diagnostics import ProviderDiagnostics

RELAY_VERSION = 'relay-failure-provenance-20261025-v1'
MODEL = 'deepseek-flash'
ENDPOINT = 'https://api.deepseek.com'
ENV_KEY = 'YOHAKU_PROVIDER_TOKEN'
MAX_BODY = 16*1024*1024


def private_path():
    return Path.home()/'.config/yohaku/codex-provider.json'


def _private_directory():
    directory = private_path().parent
    directory.mkdir(mode=0o700,parents=True,exist_ok=True)
    metadata = directory.lstat()
    if not stat.S_ISDIR(metadata.st_mode) or metadata.st_uid!=os.getuid() or metadata.st_mode & 0o077:
        raise ValueError('Yohaku provider directory must be an owned, non-symlink directory with mode 0700')
    return directory


def load():
    path = private_path()
    if not path.exists() and not path.is_symlink():
        return {'active':'openai'}
    _private_directory()
    descriptor = os.open(path,os.O_RDONLY|os.O_NOFOLLOW|os.O_CLOEXEC|os.O_NONBLOCK)
    try:
        metadata = os.fstat(descriptor)
        if not stat.S_ISREG(metadata.st_mode) or metadata.st_uid!=os.getuid() or metadata.st_mode & 0o077 or metadata.st_size>8192:
            raise ValueError('Yohaku provider config must be an owned regular file with mode 0600')
        with os.fdopen(descriptor,'r',closefd=False) as stream:
            value = json.loads(stream.read(8193))
    except (ValueError,UnicodeError):
        raise ValueError('Invalid Yohaku private provider config') from None
    finally:
        os.close(descriptor)
    if not isinstance(value,dict) or value.get('active') not in ('openai','deepseek') or set(value)-{'active','deepseek_key','openai_model'}:
        raise ValueError('Invalid Yohaku private provider config')
    if 'openai_model' in value:
        model=value['openai_model']
        if not isinstance(model,str) or not 1<=len(model)<=200 or any(c not in 'abcdefghijklmnopqrstuvwxyzABCDEFGHIJKLMNOPQRSTUVWXYZ0123456789._-:/' for c in model):
            raise ValueError('Invalid selected OpenAI model')
    if 'deepseek_key' in value:
        validate_key(value['deepseek_key'])
    if value['active']=='deepseek' and not value.get('deepseek_key'):
        raise ValueError('DeepSeek requires a key configured with the local hidden prompt')
    return value


def validate_key(key):
    if not isinstance(key,str) or not 16<=len(key)<=512 or any(not 33<=ord(character)<=126 for character in key):
        raise ValueError('Invalid API key format; no key value is displayed')


def save(value):
    directory = _private_directory()
    descriptor,temporary = tempfile.mkstemp(prefix='.provider-',dir=directory)
    try:
        with os.fdopen(descriptor,'w') as stream:
            json.dump(value,stream)
            stream.write('\n')
            stream.flush()
            os.fsync(stream.fileno())
        os.replace(temporary,private_path())
    finally:
        if os.path.exists(temporary):
            os.unlink(temporary)


def catalog():
    return {'models':[{
        'slug':MODEL,'display_name':'DeepSeek-Flash','description':'Yohaku DeepSeek Responses provider',
        'default_reasoning_level':'high',
        'supported_reasoning_levels':[{'effort':level,'description':level} for level in ('low','high','max')],
        'shell_type':'shell_command','visibility':'list','supported_in_api':True,'priority':1,
        'base_instructions':'Use only the provided Yohaku world tools. Follow the developer instructions and current world permissions.',
        'supports_reasoning_summaries':True,'default_reasoning_summary':'none',
        'support_verbosity':True,'default_verbosity':'low','apply_patch_tool_type':'freeform',
        'web_search_tool_type':'text','truncation_policy':{'mode':'tokens','limit':10000},
        'supports_parallel_tool_calls':True,'context_window':1048576,'effective_context_window_percent':95,
        'experimental_supported_tools':[],'input_modalities':['text','image'],
        'prefer_websockets':False,'supports_search_tool':False,'minimal_client_version':'0.144.0'
    }]}


def runtime_config(home,base_url,policy):
    return '''model = "deepseek-flash"
model_provider = "deepseek"
forced_login_method = "api"
model_reasoning_effort = "high"
web_search = "disabled"
model_catalog_json = '''+json.dumps(str(home/'models.json'))+'\n'+policy+'''
[model_providers.deepseek]
name = "DeepSeek"
base_url = '''+json.dumps(base_url)+'''
wire_api = "responses"
env_key = "YOHAKU_PROVIDER_TOKEN"
requires_openai_auth = false
request_max_retries = 0
stream_max_retries = 0
'''


class Relay:
    def __init__(self,key):
        validate_key(key)
        self.key = key
        self.token = secrets.token_urlsafe(32)
        self.records = []
        self.request_count = 0
        self.lock = threading.Lock()
        self.recovery=Recovery()
        self.diagnostics=ProviderDiagnostics('responses',MODEL,ENDPOINT+'/responses',key,backend='codex-relay')
        relay = self

        class Handler(BaseHTTPRequestHandler):
            protocol_version = 'HTTP/1.1'

            def log_message(self,*arguments):
                pass

            def reject(self,status,detail=None):
                payload=json.dumps({'error':detail},ensure_ascii=True).encode() if detail else b''
                self.send_response(status)
                if detail:self.send_header('Content-Type','application/json')
                self.send_header('Content-Length',str(len(payload)))
                self.send_header('Connection','close')
                self.end_headers()
                self.close_connection = True
                if payload:self.wfile.write(payload)

            def do_POST(self):
                upstream = None
                started = False
                try:
                    self.connection.settimeout(120)
                    if self.path!='/responses':
                        return self.reject(404)
                    if not hmac.compare_digest(self.headers.get('Authorization',''),'Bearer '+relay.token):
                        return self.reject(403)
                    try:length = int(self.headers.get('Content-Length','0'))
                    except ValueError:return self.reject(400)
                    if self.headers.get('Transfer-Encoding') or not 0<length<=MAX_BODY:
                        return self.reject(413)
                    body = self.rfile.read(length)
                    try:request = json.loads(body)
                    except (ValueError,UnicodeError):return self.reject(400)
                    if not isinstance(request,dict) or request.get('model')!=MODEL or not isinstance(request.get('reasoning'),dict) or request['reasoning'].get('effort')!='high':
                        return self.reject(400)
                    record = {'relay_version':RELAY_VERSION,'target':ENDPOINT+'/responses','request_model':request['model'],
                              'reasoning_effort':'high','response_models':[],'usage':None,'completed':False}
                    with relay.lock:
                        relay.request_count += 1
                        record['request_index'] = relay.request_count
                        first_sequence=relay.diagnostics.sequence
                        relay.records.append(record)
                        relay.records = relay.records[-128:]
                    # Nothing reaches Codex until transport and the terminal SSE event
                    # are complete. Failed-attempt deltas/tool arguments are discarded.
                    payload=relay.recovery.run(lambda:relay.fetch(body,request,record),
                        lambda:{'phase':relay.diagnostics.transport.failed_phase,
                                'http_status':relay.diagnostics.transport.origin_status,
                                'request_sequence':relay.diagnostics.sequence,
                                'diagnostic_path':relay.diagnostics.last_failure_path})
                    record['recovery']=list(relay.recovery.request_events)
                    record['upstream_attempt_count']=relay.diagnostics.sequence-first_sequence
                    self.send_response(200)
                    self.send_header('Content-Type','text/event-stream')
                    self.send_header('Connection','close')
                    started = True
                    record['downstream_started']=True
                    record['relay_http_status']=200
                    self.end_headers()
                    self.close_connection = True
                    record['delivery_started']=True
                    relay.recovery.check()
                    self.wfile.write(payload)
                    self.wfile.flush()
                except (OSError,ValueError,TypeError,http.client.HTTPException,RecoveryCancelled) as error:
                    if 'record' in locals():
                        record['recovery']=list(relay.recovery.request_events)
                        record['upstream_attempt_count']=relay.diagnostics.sequence-first_sequence
                        if started and not isinstance(error,RecoveryCancelled):
                            relay.diagnostics.transport.failure('downstream_delivery',error)
                            relay.diagnostics.failure(error,request,relay.diagnostics.transport.origin_status,stage='downstream_delivery')
                        failure=relay.failure(error,record)
                        record['failure']=failure
                        record['relay_http_status']=200 if started else 502
                    else:
                        return self.reject(400,{'type':'yohaku_bad_request','message':'Invalid local relay request'})
                    if not started:
                        self.reject(502,{'type':'yohaku_'+failure['kind'],'code':failure['kind'],
                            'message':relay.failure_message(failure),**failure})
                finally:
                    if upstream is not None:
                        upstream.close()

        self.server = HTTPServer(('127.0.0.1',0),Handler)
        self.server.timeout = 1
        self.thread = threading.Thread(target=self.server.serve_forever,daemon=True)
        self.thread.start()

    def failure(self, error, record):
        trace=self.diagnostics.transport
        status=trace.origin_status if trace else None
        phase=trace.failed_phase if trace else None
        if record.get('upstream_attempt_count')==0:status=None;phase=None
        if isinstance(error,RecoveryCancelled):kind='cancelled';phase='cancelled'
        elif record.get('downstream_started') or record.get('delivery_started'):kind='downstream_delivery';phase='downstream_delivery'
        elif isinstance(error,urllib.error.HTTPError):kind='upstream_http'
        elif isinstance(error,(ValueError,TypeError)):kind='response_parse' if status is not None else 'local_adapter'
        else:kind='connection' if status is None else 'response_read'
        return {'kind':kind,'upstream_http_status':status,
            'phase':phase or ('response_decode' if kind=='response_parse' else kind),
            'exception_type':type(error).__name__,'diagnostic_id':self.diagnostics.session+'-'+str(self.diagnostics.sequence),
            'diagnostic_path':self.diagnostics.last_failure_path,
            'attempts':record.get('upstream_attempt_count',0),'retryable':retryable(error) and not (record.get('downstream_started') or record.get('delivery_started')),
            'retry_exhausted':bool(self.recovery.request_events and self.recovery.request_events[-1]['attempt']>=self.recovery.max_retries+1 and retryable(error)),
            'tools_delivered':'unknown' if record.get('delivery_started') else False}

    @staticmethod
    def failure_message(failure):
        status=failure['upstream_http_status']
        return ('Yohaku relay failure: kind='+failure['kind']+'; upstream HTTP '+
            ('not observed' if status is None else str(status))+'; phase='+failure['phase']+
            '; attempts='+str(failure['attempts'])+'; diagnostic='+failure['diagnostic_id']+
            ('; no incomplete model/tool output delivered' if failure.get('tools_delivered') is False else '; downstream delivery outcome unknown; inspect before replay'))

    def configure_recovery(self,config,emit=lambda event:None):
        # Configured once before turn/start. No change to account or native routes.
        self.recovery=Recovery(config)
        self.recovery.notify=emit

    def sample_response(self,response):
        # Bounded error evidence only, including 200 responses with a wrong MIME type.
        if not hasattr(response,'read'):return
        try:
            sample=response.read(8193)
            remaining=getattr(response,'length',None)
            self.diagnostics.observe_body(sample[:8192],len(sample)<=8192 and not (type(remaining) is int and remaining>0),len(sample)>8192)
        except (OSError,http.client.HTTPException) as error:
            self.diagnostics.observe_body(getattr(error,'partial',b'')[:8192],False)

    def fetch(self,body,request,record):
        self.diagnostics.begin(timeout=120,request_bytes=len(body),messages=request.get('input'),recovery=self.recovery)
        trace=self.diagnostics.transport
        trace.metadata['connection_policy']='http.client HTTPS; new connection per upstream attempt; closed in finally'
        record['confirmed_tool_results']=trace.metadata.get('confirmed_tool_results',0)
        record['request_bytes']=len(body)
        upstream=None;response=None;parts=[];size=0;read_complete=False;decode_started=False
        try:
            upstream=ObservedHTTPS(trace,'api.deepseek.com',timeout=120)
            upstream.request('POST','/responses',body=body,headers={
                'Authorization':'Bearer '+self.key,'Content-Type':'application/json','Accept':'text/event-stream',
                'Accept-Encoding':'identity','User-Agent':'Yohaku-Codex-DeepSeek/1'})
            response=upstream.getresponse()
            record['http_status']=response.status
            trace.origin_status=response.status
            self.diagnostics.observe_response(response)
            if response.status!=200:
                import urllib.error
                self.sample_response(response)
                error=urllib.error.HTTPError(ENDPOINT+'/responses',response.status,'Provider status',{},None)
                error.close()
                raise error
            if 'text/event-stream' not in response.getheader('Content-Type','').lower():
                self.sample_response(response)
                raise ValueError('Provider did not return an event stream')
            trace.event('response_read_start')
            try:
                while True:
                    self.recovery.check()
                    line=response.readline(MAX_BODY+1)
                    if not line:break
                    size+=len(line)
                    if size>MAX_BODY:raise ValueError('Provider stream exceeds 16 MiB buffered limit')
                    parts.append(line)
                if type(getattr(response,'length',None)) is int and response.length>0:raise http.client.IncompleteRead(b'',response.length)
            except (OSError,http.client.HTTPException) as error:
                trace.failure('response_read',error);raise
            read_complete=True
            self.diagnostics.observe_body(b''.join(parts),True)
            trace.event('response_read_complete',bytes_read=size)
            payload=b''.join(parts)
            # SSE supports multiple data lines in a frame. JSON is validated only
            # after receipt; partial function-call fragments never reach Codex.
            decode_started=True
            trace.event('response_decode_start')
            completed=False
            for frame_index,frame in enumerate(payload.replace(b'\r\n',b'\n').split(b'\n\n')):
                data=b'\n'.join(line[5:].lstrip(b' ') for line in frame.split(b'\n') if line.startswith(b'data:'))
                if not data or data==b'[DONE]':continue
                try:event=json.loads(data)
                except (json.JSONDecodeError,UnicodeDecodeError) as error:
                    self.diagnostics.observe_parse_error(payload,frame_index,data,error)
                    raise
                if not isinstance(event,dict):raise ValueError('Invalid provider event')
                if event.get('type') in ('response.failed','response.incomplete','error'):
                    raise ValueError('Provider reported an incomplete or failed response; no tools delivered')
                if event.get('type')=='response.completed':
                    if not isinstance(event.get('response'),dict) or event['response'].get('status','completed')!='completed':
                        raise ValueError('Invalid terminal response status')
                    completed=True
            self.diagnostics.response_evidence['terminal_event_received']=completed
            if not completed:
                trace.failure('response_read',http.client.IncompleteRead(b''))
                raise http.client.IncompleteRead(b'')
            for line in payload.splitlines():self.observe(line,record)
            self.diagnostics.succeeded(response.status)
            return payload
        except (OSError,ValueError,TypeError,http.client.HTTPException) as error:
            if not trace.failed_phase:
                trace.failure('response_decode' if decode_started or isinstance(error,ValueError) else 'http_response' if trace.origin_status is not None else 'connection',error)
            if parts and not read_complete:
                received=b''.join(parts)+getattr(error,'partial',b'')
                self.diagnostics.observe_body(received[:MAX_BODY],False,len(received)>MAX_BODY)
            self.diagnostics.failure(error,request,trace.origin_status)
            raise
        finally:
            if response is not None:response.close()
            if upstream is not None:upstream.close()

    @property
    def base_url(self):
        return 'http://127.0.0.1:'+str(self.server.server_port)

    def observe(self,line,record):
        if not line.startswith(b'data:'):
            return
        try:
            event = json.loads(line[5:])
            response = event.get('response',{})
            model = response.get('model')
            if isinstance(model,str) and len(model)<=128 and model not in record['response_models'] and len(record['response_models'])<8:
                record['response_models'].append(model)
            if event.get('type')=='response.completed':
                usage = response.get('usage',{})
                record['usage'] = {name:value for name,value in usage.items() if name in (
                    'input_tokens','output_tokens','total_tokens') and type(value) is int and value>=0}
                record['completed'] = True
        except (ValueError,TypeError,AttributeError):
            pass

    def evidence(self):
        with self.lock:
            return json.loads(json.dumps(self.records))

    def close(self):
        self.recovery.finish_budget()
        self.recovery.cancel()
        self.server.shutdown()
        self.server.server_close()
        self.thread.join(timeout=5)
        self.key = ''
        self.token = ''


def main():
    parser = argparse.ArgumentParser(description=__doc__)
    parser.add_argument('mode',choices=['configure-deepseek','use-deepseek','use-openai','status'])
    arguments = parser.parse_args()
    try:
        value = load()
        if arguments.mode=='configure-deepseek':
            if not sys.stdin.isatty():
                raise ValueError('Use a local interactive terminal for hidden key entry; do not pipe a key')
            with warnings.catch_warnings():
                warnings.simplefilter('error',getpass.GetPassWarning)
                key = getpass.getpass('DeepSeek API key (hidden; Yohaku only): ')
            validate_key(key)
            value = {'active':'deepseek','deepseek_key':key}
            save(value)
        elif arguments.mode.startswith('use-'):
            value['active'] = arguments.mode[4:]
            if value['active']=='deepseek' and not value.get('deepseek_key'):
                raise ValueError('First run configure-deepseek in a local terminal')
            save(value)
        print(json.dumps({'active':value['active'],'key_configured':bool(value.get('deepseek_key')),
                          'model':MODEL if value['active']=='deepseek' else value.get('openai_model','Yohaku turn config model'),
                          'private_path':str(private_path())}))
    except (OSError,ValueError,EOFError,getpass.GetPassWarning):
        print('Provider configuration failed; check private file permissions and use a local hidden prompt. No key is displayed.',file=sys.stderr)
        return 1
    return 0


if __name__=='__main__':
    sys.exit(main())
````

## 输出与关联
```yaml
output: "agent/yohaku_codex_provider.py"
```

## AI 笔记
```yaml
{"source_version": 1, "writer": "Yohaku release packaging", "summary": "Portable release source. Scope and packaging adjustments are documented in PROVENANCE.md."}
```
